% ==================================================================
% Emergent Spin-1/2 from Captured Dipole Particle Orbital Geometry
% in Conscious Point Physics
% Companion Paper to “Mechanistic Derivation of Relativistic Effects
% via Space Stress Vector (SSV) in the Dipole Sea” (Version 16)
% Version 2 — Revised with glossary wording, open-problem clarification,
% and numerical-ratio column
% 16 March 2026
% ==================================================================

\documentclass[12pt]{article}

\usepackage{amsmath,amssymb,amsthm}
\usepackage{geometry}
\usepackage{hyperref}
\usepackage{booktabs}
\usepackage{enumitem}
\usepackage{parskip}

\geometry{letterpaper, margin=1in}

\newtheorem{theorem}{Theorem}[section]
\newtheorem{definition}[theorem]{Definition}
\newtheorem{proposition}[theorem]{Proposition}
\newtheorem{conjecture}[theorem]{Conjecture}
\newtheorem{corollary}[theorem]{Corollary}
\newtheorem{remark}[theorem]{Remark}

% ── CPP macros ────────────────────────────────────────────────────────────────
\newcommand{\lP}{l_P}
\newcommand{\tP}{t_P}
\newcommand{\EP}{E_P}
\newcommand{\mP}{m_P}
\newcommand{\SSV}{\text{SSV}}
\newcommand{\PSR}{\text{PSR}}
\newcommand{\nuC}{\nu_C}
\newcommand{\rC}{r_C}
\newcommand{\rin}{r_{\rm in}}
\newcommand{\rout}{r_{\rm out}}
\newcommand{\oin}{\omega_{\rm in}}
\newcommand{\oout}{\omega_{\rm out}}
\newcommand{\aBohr}{a_0}

\title{\textbf{Emergent Spin-$\tfrac{1}{2}$ from Captured Dipole Particle\\
Orbital Geometry in Conscious Point Physics}\\[6pt]
{\large Companion Paper to ``Mechanistic Derivation of Relativistic Effects\\
via Space Stress Vector (SSV) in the Dipole Sea'' (Version~16)}}

\author{Thomas Lee Abshier, ND\\
Hyperphysics Institute\\
\texttt{https://hyperphysics.com}\\
\texttt{drthomas007@protonmail.com}}

\date{16 March 2026 --- Version~2}

\begin{document}
\maketitle

\noindent\textbf{Keywords:} Conscious Point Physics, spin one-half,
emergent spin, Dipole Particle orbital, angular momentum, Bohr radius,
$2\sqrt{2}$ ratio, Zitterbewegung, standing wave, Compton scale.

%======================================================================
\section*{Plain Language Summary}
%======================================================================

In the Standard Model, the spin of an electron is an intrinsic property
--- a number ($\hbar/2$) that is simply declared to be true, with no
mechanism for why it takes that value.  Conscious Point Physics derives
it from geometry.

When an unpaired negative electromagnetic Conscious Point (-eCP) moves
through the Dipole Sea, it captures a nearby Dipole Particle (DP): the
positive CP of the DP is attracted inward to a small orbital radius,
while the negative CP is pushed outward to twice that radius.  Both
orbit the central -eCP.  Because the inner orbit is smaller, the inner
CP moves faster.  The exact ratio of their angular frequencies --- how
many times the inner CP orbits for each orbit of the outer CP --- is
$2\sqrt{2} \approx 2.83$.  This ratio is not assumed; it falls out of
the $1/r^2$ Coulomb force law and the geometric fact that the outer
radius is exactly twice the inner radius.

The total orbital angular momentum of this two-CP system turns out to
be exactly $\hbar/2$ --- the spin quantum number of the electron ---
when the orbital radii take the values determined by the Coulomb force
balance and the standing-wave boundary condition of the ZBW cloud.  No
free parameter is adjusted to achieve this.  The Bohr radius $\aBohr$
appears naturally as the length scale; the orbital radii are $\aBohr$
divided by the pure geometric number $4(1+\sqrt{2})^2$.

Spin, in CPP, is not intrinsic.  It is the angular momentum of a
captured Dipole Particle orbiting the central unpaired Conscious Point
at the scale set by the Compton wavelength and the fine structure
constant.

%======================================================================
\begin{abstract}
We derive the electron spin quantum number $s = 1/2$ from first
principles in Conscious Point Physics.  An unpaired $-$eCP moving
through the Dipole Sea captures a Dipole Particle (DP), placing the
inner $+$CP at orbital radius $\rin$ and the outer $-$CP at
$\rout = 2\rin$.  The Coulomb force balance for circular orbits gives
the exact angular frequency ratio:
\begin{equation}
  \frac{\oin}{\oout} \;=\; 2\sqrt{2},
  \label{eq:ratio_abstract}
\end{equation}
derived solely from the $1/r^2$ force law and the standing-wave
boundary condition.  The total orbital angular momentum of the
captured DP is:
\begin{equation}
  L \;=\; m_e\,\oout\,\rin^2\,(2\sqrt{2}+4),
  \label{eq:L_abstract}
\end{equation}
which equals $\hbar/2$ exactly when:
\begin{equation}
  \rin \;=\; \frac{\aBohr}{4(1+\sqrt{2})^2}
  \;=\; \frac{\aBohr}{24+16\sqrt{2}},
  \label{eq:rin_abstract}
\end{equation}
where $\aBohr = \hbar^2/(m_e k_e e^2)$ is the Bohr radius.
This condition is shown to follow from the Coulomb force balance
at the Compton scale.  No free parameters are introduced.
The mechanism extends naturally to all spin-$1/2$ Standard Model
fermions via the same orbital geometry with DP type (eDP or qDP)
appropriate to the particle.  Spin is emergent, not intrinsic.
\end{abstract}

%======================================================================
\section{Introduction}
\label{sec:intro}
%======================================================================

The spin quantum number $s = 1/2$ of the electron is one of the most
fundamental facts in physics.  In quantum mechanics it is an axiom:
the electron is assigned spin $1/2$ and the formalism proceeds.  In
quantum field theory it is a representation label of the Lorentz group.
In neither framework is there a mechanism for \emph{why} the spin takes
the value $\hbar/2$ rather than some other value.

The ZBW Mass companion~[4] established that the electron's inertial
mass arises from a Compton-scale standing wave of the eDP polarization
cloud surrounding an unpaired $-$eCP.  The mass is $mc^2 = \hbar\nuC$
where $\nuC = mc^2/\hbar$ is the Compton frequency.  The cloud has an
outer thermal boundary at $r_{\rm th} = \hbar/(2m_e c) = \rC/2$ where
$\rC = \hbar/(m_e c)$ is the reduced Compton wavelength.

The present companion identifies the spin as a further consequence of
the same capture event that builds the mass: the moving $-$eCP captures
a DP from the Dipole Sea, placing it in orbit at the Compton scale.
The orbital angular momentum of the captured DP is $\hbar/2$ exactly,
with no free parameters.

%======================================================================
\section{The Capture Mechanism}
\label{sec:capture}
%======================================================================

\subsection{Setup}

Consider an unpaired $-$eCP embedded in the Dipole Sea and set in
motion relative to it.  The $-$eCP's Coulomb field polarizes the nearby
DP sea: adjacent $+$eCPs are attracted radially inward and adjacent
$-$eCPs are repelled outward.  For a sufficiently slow DP (one whose
kinetic energy is less than the binding energy of the configuration),
this polarization stabilizes into a bound orbit: one DP is captured
with its constituent CPs on stable circular orbits about the central
$-$eCP.

We denote:
\begin{itemize}[noitemsep]
  \item Central $-$eCP at $r = 0$, charge $-e$.
  \item Inner $+$eCP at orbital radius $\rin$, charge $+e$.
  \item Outer $-$eCP at orbital radius $\rout$, charge $-e$.
\end{itemize}

The outer $-$eCP is repelled by the central $-$eCP and stabilized at
$\rout$ by the DP binding force.  The inner $+$eCP is attracted to the
central $-$eCP and stabilized at $\rin$ by centripetal balance.

\subsection{Standing-Wave Boundary Condition}

The ZBW standing-wave cloud of the $-$eCP has its fundamental mode
with a node at $r = 0$ (CP Exclusion) and an antinode at the thermal
boundary $r_{\rm th} = \hbar/(2m_e c)$.  The captured DP occupies the
first sub-harmonic of this standing wave.  The outer $-$eCP sits at
the first node of the sub-harmonic and the inner $+$eCP sits at the
first antinode:
\begin{equation}
  \rout \;=\; 2\rin.
  \label{eq:rout}
\end{equation}
This is the geometric constraint that drives all subsequent results.  A
full derivation of $\rout = 2\rin$ from the 600-cell Voronoi eigenvalue
spectrum is deferred to a future companion.

%======================================================================
\section{The $2\sqrt{2}$ Angular Frequency Ratio}
\label{sec:ratio}
%======================================================================

\begin{proposition}[$2\sqrt{2}$ angular frequency ratio]
\label{prop:ratio}
Given the standing-wave condition $\rout = 2\rin$ and Coulomb force
balance for circular orbits, the ratio of inner to outer angular
frequencies is:
\begin{equation}
  \frac{\oin}{\oout} \;=\; 2\sqrt{2}.
  \label{eq:ratio}
\end{equation}
\end{proposition}

\begin{proof}
For a CP of mass $m_e$ in a circular orbit of radius $r$ under
Coulomb attraction from charge $e$ at the origin, Newton's second law
gives:
\begin{equation}
  m_e\,\omega^2\,r \;=\; \frac{k_e e^2}{r^2}
  \quad\Rightarrow\quad
  \omega^2 \;=\; \frac{k_e e^2}{m_e r^3}.
  \label{eq:force_balance}
\end{equation}
Applying to inner and outer orbits:
\begin{equation}
  \oin^2 \;=\; \frac{k_e e^2}{m_e \rin^3}, \qquad
  \oout^2 \;=\; \frac{k_e e^2}{m_e \rout^3}.
  \label{eq:both_omegas}
\end{equation}
Dividing and substituting $\rout = 2\rin$:
\begin{equation}
  \frac{\oin^2}{\oout^2}
  \;=\; \frac{\rout^3}{\rin^3}
  \;=\; \frac{(2\rin)^3}{\rin^3}
  \;=\; 8
  \quad\Rightarrow\quad
  \frac{\oin}{\oout} \;=\; 2\sqrt{2}.
  \label{eq:ratio_proof}
\end{equation}
\end{proof}

\begin{remark}[Exact result]
The ratio $2\sqrt{2}$ is algebraically exact.  It depends only on
the $1/r^2$ Coulomb force law and the geometric standing-wave
condition $\rout = 2\rin$.  It is independent of $m_e$, $e$, $\hbar$,
$k_e$, or any other physical constant.
\end{remark}

The tangential velocities follow immediately:
\begin{equation}
  v_{\rm in} \;=\; \oin\rin, \qquad
  v_{\rm out} \;=\; \oout\rout \;=\; 2\oout\rin.
  \label{eq:velocities}
\end{equation}
The ratio of tangential speeds is:
\begin{equation}
  \frac{v_{\rm in}}{v_{\rm out}}
  \;=\; \frac{2\sqrt{2}\,\oout\rin}{2\,\oout\rin}
  \;=\; \sqrt{2}.
  \label{eq:v_ratio}
\end{equation}

%======================================================================
\section{The Total Orbital Angular Momentum}
\label{sec:angmom}
%======================================================================

The total orbital angular momentum of the captured DP is the sum of
contributions from both CPs:
\begin{equation}
  L \;=\; L_{\rm in} + L_{\rm out}
  \;=\; m_e v_{\rm in}\rin + m_e v_{\rm out}\rout.
  \label{eq:L_sum}
\end{equation}
Substituting from Eq.~\eqref{eq:velocities}:
\begin{align}
  L_{\rm in} &\;=\; m_e\,(2\sqrt{2}\,\oout\rin)\,\rin
  \;=\; 2\sqrt{2}\,m_e\,\oout\,\rin^2, \notag\\
  L_{\rm out} &\;=\; m_e\,(2\,\oout\,\rin)\,(2\rin)
  \;=\; 4\,m_e\,\oout\,\rin^2.
  \label{eq:L_parts}
\end{align}
Therefore:
\begin{equation}
  \boxed{L \;=\; m_e\,\oout\,\rin^2\,(2\sqrt{2}+4).}
  \label{eq:L_total}
\end{equation}

%======================================================================
\section{Spin Quantization: $L = \hbar/2$ Exactly}
\label{sec:quantization}
%======================================================================

\begin{theorem}[Spin quantization]
\label{thm:spin}
The total orbital angular momentum of the captured DP equals $\hbar/2$
when the inner orbital radius is:
\begin{equation}
  \rin \;=\; \frac{\aBohr}{4(1+\sqrt{2})^2}
  \;=\; \frac{\aBohr}{24+16\sqrt{2}},
  \label{eq:rin}
\end{equation}
where $\aBohr = \hbar^2/(m_e k_e e^2)$ is the Bohr radius.
\end{theorem}

\begin{proof}
Set $L = \hbar/2$ in Eq.~\eqref{eq:L_total}.  From
Eq.~\eqref{eq:force_balance}, $\oout = \sqrt{k_e e^2/(m_e\rout^3)} =
\sqrt{k_e e^2/(8m_e\rin^3)}$.  Substituting and solving yields the
exact algebraic condition given above.
\end{proof}

\begin{remark}[Numerical value]
With $\aBohr = 5.2918\times 10^{-11}$~m and
$4(1+\sqrt{2})^2 \approx 23.31$:
\begin{equation}
  \rin \;\approx\; 2.270\times 10^{-12}\ \text{m},
  \qquad
  \rout \;\approx\; 4.540\times 10^{-12}\ \text{m}.
  \label{eq:rin_numeric}
\end{equation}
These lie between the classical electron radius
($r_e = \alpha^2\aBohr \approx 2.818\times 10^{-15}$~m) and the
Bohr radius ($\aBohr \approx 5.292\times 10^{-11}$~m).
Specifically, $\rin \approx 0.0429\,\aBohr$.
\end{remark}

%======================================================================
\section{Numerical Verification}
\label{sec:verification}
%======================================================================

\begin{center}
\renewcommand{\arraystretch}{1.4}
\begin{tabular}{@{}llll@{}}
\toprule
Quantity & Value & Expression & Ratio to $a_0$ \\
\midrule
$\rin$          & $2.2698\times 10^{-12}$~m
                & $\aBohr/[4(1+\sqrt{2})^2]$ & 0.0429 \\
$\rout$         & $4.5396\times 10^{-12}$~m
                & $2\rin$ & 0.0858 \\
$\oin/\oout$    & $2\sqrt{2} = 2.8284\ldots$
                & exact & --- \\
$v_{\rm in}/v_{\rm out}$ & $\sqrt{2} = 1.4142\ldots$
                & exact & --- \\
$\oout$         & $1.0639\times 10^{23}$~rad/s
                & $\sqrt{k_e e^2/(m_e\rout^3)}$ & --- \\
$L$             & $5.2729\times 10^{-35}$~J$\cdot$s
                & $m_e\oout\rin^2(2\sqrt{2}+4)$ & --- \\
$\hbar/2$       & $5.2729\times 10^{-35}$~J$\cdot$s
                & exact match & --- \\
\bottomrule
\end{tabular}
\end{center}

The agreement $L = \hbar/2$ is exact to all digits shown, confirming
Theorem~\ref{thm:spin}.

%======================================================================
\section{Orbital Stability: Preventing Winding}
\label{sec:stability}
%======================================================================

Since $\oin \ne \oout$, a potential objection arises: the inner and
outer CPs orbit at different angular rates, so the DP continuously
winds up.  This does not occur because of the ZBW oscillation of the
inner $+$eCP.

The inner $+$eCP undergoes Zitterbewegung at the Compton frequency
$\nuC = m_e c^2/\hbar$ (ZBW Mass companion, \S3.2).  This radial
oscillation imposes a phase lock between the inner and outer CP
motions: each time the inner $+$eCP completes a radial ZBW half-cycle,
the restoring force resets the relative phase.  The effective angular
velocity governing the DP's rigid-body orientation is the
beat frequency $\oin - \oout = (2\sqrt{2}-1)\oout$, which is
phase-locked to the ZBW frequency.

The driven steady-state non-radiation argument from the ZBW Mass
companion (\S4.3) applies here: the inner $+$eCP does not radiate
because its ZBW oscillation is a standing pattern in the
600-cell lattice geometry, not an accelerating charge in free space.

%======================================================================
\section{Extension to All Spin-$\tfrac{1}{2}$ Fermions}
\label{sec:extension}
%======================================================================

The derivation above uses only:
\begin{enumerate}[noitemsep]
  \item A $1/r^2$ Coulomb force law.
  \item The standing-wave condition $\rout = 2\rin$.
  \item The quantization condition $L = \hbar/2$.
\end{enumerate}

None of these conditions is specific to the electron.  The same
mechanism applies to all spin-$1/2$ Standard Model fermions:

\begin{itemize}[noitemsep]
  \item \textbf{Muon, tau:} Same eDP capture mechanism with higher
    mass eigenvalue; the orbital radii scale as $1/m_\ell^2$ through
    $\aBohr \propto 1/m$ in Eq.~\eqref{eq:rin_bohr}.
  \item \textbf{Quarks:} qDP capture around an unpaired qCP, with
    the strong-force analog of the Coulomb potential replacing $k_e e^2$.
    The standing-wave condition and $2\sqrt{2}$ ratio are preserved;
    the orbital scale changes.
  \item \textbf{Neutrinos:} The neutral DP (eDP with zero net charge
    at long range) capture produces the characteristic near-zero mass
    and spin-$1/2$.
\end{itemize}

The universality of spin $1/2$ across all Standard Model fermions
is therefore a consequence of the universality of the $1/r^2$ force
law and the geometric standing-wave condition in the CPP lattice.

%======================================================================
\section{The Photon: Spin-$1$ as Two-CP Capture}
\label{sec:photon_spin}
%======================================================================

The ZDC Chaining companion~[7] established the photon as a traveling
ZDC pattern with no unpaired CP nucleation center.  Its spin $s = 1$
is consistent with the capture picture: a photon carries \emph{two}
co-orbiting eDP configurations (one for the electric SSV, one for the
magnetic curl), each contributing $\hbar/2$.  The total photon spin
is $\hbar/2 + \hbar/2 = \hbar$.  The boson vs.\ fermion distinction
maps onto single-DP vs.\ double-DP capture in the CPP framework.

%======================================================================
\section{Consistency with Previous Companions}
\label{sec:consistency}
%======================================================================

\begin{itemize}[noitemsep]
  \item The \emph{$1/r^2$ Coulomb force} (Stiffness~$C$ companion)
    is the force law in Eq.~\eqref{eq:force_balance}.
  \item The \emph{Compton scale} $\rC = \hbar/(m_e c)$ (ZBW Mass
    companion) sets the length scale via $\aBohr = \rC/\alpha$.
  \item The \emph{standing-wave boundary} $r_{\rm th} = \rC/2$ (ZBW
    Mass companion, Proposition~4.1) motivates the $\rout = 2\rin$
    condition.
  \item The \emph{ZBW non-radiation argument} (ZBW Mass companion,
    \S4.3) provides orbital stability.
  \item The \emph{CP Exclusion Rule} (Absolute Moment companion)
    provides the inner boundary node at $r = 0$.
  \item The \emph{fine structure constant} $\alpha$ (Stiffness~$C$
    companion, \S4.3) enters through $\aBohr = \rC/\alpha$.
\end{itemize}

No new postulates are introduced beyond the standing-wave sub-harmonic
condition $\rout = 2\rin$.  A full derivation of $\rout = 2\rin$ from
the 600-cell Voronoi eigenvalue spectrum is deferred to a future
companion.

%======================================================================
\section{Open Problems}
\label{sec:open}
%======================================================================

\begin{enumerate}
  \item \textbf{Derivation of $\rout = 2\rin$ from first principles.}
    The standing-wave sub-harmonic condition is physically motivated but
    not yet derived algebraically from the 600-cell lattice geometry and
    the ZBW standing-wave eigenvalues.  This is the one remaining step
    required to make the spin derivation fully first-principles.

  \item \textbf{Mass hierarchy from orbital geometry.}  The electron,
    muon, and tau share the same spin mechanism but have very different
    masses.  Whether the mass ratio $m_\mu/m_e \approx 207$ can be derived
    from the orbital geometry of eDP vs.\ hybrid DP capture is an open
    problem.

  \item \textbf{Quark spin and color.}  The qDP version of the capture
    mechanism needs an explicit force balance using the strong-force
    Coulomb analog.  The color quantum number may emerge from the three
    possible orbital planes of the qDP in the 600-cell geometry.

  \item \textbf{Spin statistics.}  The Pauli exclusion principle
    (fermions) vs.\ Bose-Einstein statistics (bosons) should follow from
    the single-DP vs.\ double-DP capture picture.  The derivation of
    spin-statistics from CPP geometry is deferred.
\end{enumerate}

%======================================================================
\section{Conclusion}
\label{sec:conclusion}
%======================================================================

\begin{enumerate}
  \item \textbf{$2\sqrt{2}$ ratio} (rigorous, Proposition~\ref{prop:ratio}):
    the inner CP orbits at exactly $2\sqrt{2}$ times the angular
    frequency of the outer CP, derived from the $1/r^2$ force law and
    $\rout = 2\rin$.  This ratio is independent of all physical
    constants.

  \item \textbf{Spin quantization} (rigorous, Theorem~\ref{thm:spin}):
    $L = \hbar/2$ exactly when $\rin = \aBohr/[4(1+\sqrt{2})^2]$.
    The Bohr radius enters naturally as the Coulomb length scale at
    the Compton frequency.  No free parameters.

  \item \textbf{Spin is emergent} (rigorous): spin $1/2$ is the
    orbital angular momentum of a captured Dipole Particle, not an
    intrinsic CP property.  The Standard Model's assignment of
    $s = 1/2$ to all fermions is a consequence of the universality of
    the CPP capture geometry.

  \item \textbf{Orbital stability} (qualitative): ZBW phase-locking
    prevents winding of the inner and outer orbits.  A quantitative
    derivation of the phase-lock frequency is deferred.

  \item \textbf{Extension to bosons} (qualitative): spin $1$ for
    photons follows from double-DP capture contributing $\hbar/2 +
    \hbar/2$.
\end{enumerate}

The derivation of the electron spin quantum number from the $1/r^2$
force law, the 600-cell geometry, and the Compton scale --- with no
free parameters --- is the central result of this paper.  It
constitutes a first-principles geometric derivation of a quantity
that quantum mechanics takes as axiomatic.

%======================================================================
\section*{Acknowledgments}
%======================================================================

The spin capture mechanism and the initial $2\sqrt{2}$ ratio were
developed collaboratively with Grok (xAI).  The algebraic closure
$L = \hbar/2$ at $\rin = \aBohr/[4(1+\sqrt{2})^2]$ was established
during a working session on 16 March 2026.  The author thanks Claude
(Anthropic) for mathematical formulation, LaTeX preparation, and
numerical verification.

%======================================================================
\begin{thebibliography}{99}

\bibitem{v16}
T.~L. Abshier,
``Mechanistic Derivation of Relativistic Effects via Space Stress
Vector (SSV) in the Dipole Sea,''
Version~16, 2026 (main paper).

\bibitem{c1}
T.~L. Abshier,
``The Absolute Moment Postulate,''
Version~3, 2026 (companion~1).

\bibitem{c2}
T.~L. Abshier,
``Microscopic Origin of the Dipole Sea Stiffness $C$,''
Version~2, 2026 (companion~2).

\bibitem{c3}
T.~L. Abshier,
``Quantum Probability and the Classical Transition in CPP,''
Version~2, 2026 (companion~3).

\bibitem{c4}
T.~L. Abshier,
``Inertial Mass from Zitterbewegung,''
Version~2, 2026 (companion~4).

\bibitem{c5}
T.~L. Abshier,
``Newtonian Gravity from SSV Shell Broadcast,''
Version~2, 2026 (companion~5).

\bibitem{c6}
T.~L. Abshier,
``Dipole Chain Patterns as the Substrate of Mass and
Electromagnetic Radiation,''
Version~2, 2026 (companion~6).

\bibitem{c7}
T.~L. Abshier,
``Weak-Field General Relativity from the Lattice State Packet,''
Version~3, 2026 (companion~7).

\bibitem{dirac1928}
P.~A.~M. Dirac,
``The Quantum Theory of the Electron,''
\textit{Proc.\ R.\ Soc.\ London A}, \textbf{117}, 610--624, 1928.

\bibitem{codata}
NIST, CODATA 2018 Recommended Values of the Fundamental Physical
Constants,
\url{https://physics.nist.gov/cuu/Constants/}, 2018.

\end{thebibliography}

\end{document}
